\chapter{Results}
\label{ch:results}

\section{Chapter Overview}
\label{sec:test_overview}
The evaluation includes hand tracking, touch detection for PyTorch LSTM, homography estimation, CPU latency, and macro command execution on a paper-based virtual macro keyboard \cite{ReviewVirtualKeyboard2020}. With no physical switches or electronic wiring in the paper, interaction requires 21-joint skeletal extraction, scale-normalized kinematics, deep sequence modeling, and planar homography \cite{Zhang2020MediaPipe}.

We evaluated three key engineering challenges: detection of physical surface contact without depth sensors, scale invariance across hand sizes and camera distances, and real-time execution at 12~FPS on commodity CPUs with no GPU. Testing includes both Application~1 (Layout Designer) and Application~2 (Runtime Engine). Functional and non-functional test cases are defined in Section~\ref{sec:test_plan_cases}. Section~\ref{sec:test_workflow} outlines the evaluation procedure. Quantitative benchmarks spanning model architectures, landmark subsets, window configurations, scale normalization, CPU latency, AprilTag tilt resilience, and macro action dispatch are presented in Section~\ref{sec:test_results_review}. Section~\ref{sec:test_summary} concludes the chapter.

\section{Test Plan and Structured Test Cases}
\label{sec:test_plan_cases}
A comprehensive test plan was executed to verify that the functionality of the two applications (Application~1 Layout Designer and Application~2 Runtime Engine) worked correctly, as well as checking their execution performance, stability, and reliability.

\subsection{Functional Testing Suite}
\label{subsec:test_functional}
We performed functionality tests on layout designer operations (saving and loading XML layout files, mapping action profile types), AprilTag tracking loops, deep learning touch predictions, and executing keystrokes or native system shell commands.

Table~\ref{tab:functional_test_cases} lists the ten functional test cases evaluated across the system.

\begin{table}[htbp]
\centering
\caption{Structured functional test cases and validation results.}
\label{tab:functional_test_cases}
\begin{tabular}{p{1.2cm} p{3.5cm} p{4.5cm} p{4.0cm} c}
\toprule
\textbf{Test ID} & \textbf{Functional Feature} & \textbf{Input Test Condition} & \textbf{Expected System Output} & \textbf{Status} \\
\midrule
\textbf{FT-01} & Key Placement & User adds and positions keys on A4 canvas & Key bounding boxes registered accurately & \textbf{PASS} \\
\textbf{FT-02} & AprilTag Anchors & Automatic border anchor placement & Four corner AprilTags positioned at margins & \textbf{PASS} \\
\textbf{FT-03} & PDF Export & Vector PDF rendering & High-resolution PDF with embedded tags & \textbf{PASS} \\
\textbf{FT-04} & XML Export & Structural XML file generation & XML schema matching canvas coordinates & \textbf{PASS} \\
\textbf{FT-05} & Layout Loading & XML import in Application 2 & Lookup coordinates constructed in RAM & \textbf{PASS} \\
\textbf{FT-06} & Action Profile Multiplexing & Bind keys to different action combinations and switch profiles on same layout & Profiles saved into XML layout; distinct actions triggered on identical paper sheet & \textbf{PASS} \\
\textbf{FT-07} & Tag Detection & AprilTag quad detection on paper sheet & Corner coordinates localized with sub-pixel fit & \textbf{PASS} \\
\textbf{FT-08} & SVD Homography & Compute $3 \times 3$ matrix $H$ & Planar projection corrected for camera tilt & \textbf{PASS} \\
\textbf{FT-09} & Touch Detection & 5-frame window evaluated by LSTM & Touch probability classified for active digit & \textbf{PASS} \\
\textbf{FT-10} & Action Dispatch & Mapped key press or command trigger & OS keystroke or script executed reliably & \textbf{PASS} \\
\bottomrule
\end{tabular}%
\end{table}

All ten functional test experiments succeeded with a $100\%$ pass rate. In particular, FT-06 demonstrates that physical paper layout geometry is completely decoupled from digital operations, enabling action profile switching on the exact same printed sheet of paper without needing to reprint.

\subsection{Non-Functional Testing Suite}
\label{subsec:test_non_functional}
Non-functional testing assessed system performance metrics, including processing latency, touch classification F1-score, distance scale invariance, AprilTag tilt tolerance, RAM usage, and CPU utilization during background execution.

Table~\ref{tab:non_functional_test_cases} summarizes the seven non-functional test cases evaluated on a standard desktop computer.

\begin{table}[htbp]
\centering
\caption{Structured non-functional test cases and empirical validation results.}
\label{tab:non_functional_test_cases}
\begin{tabular}{p{1.4cm} p{3.2cm} p{4.0cm} p{4.5cm} c}
\toprule
\textbf{Test ID} & \textbf{Quality Attribute} & \textbf{Target Performance Metric} & \textbf{Measured Empirical Output} & \textbf{Status} \\
\midrule
\textbf{NFT-01} & Real-Time Pipeline Latency & End-to-end frame latency $\le 83.3\text{ ms}$ & \textbf{29.09 ms} ($\approx 34.4\text{ FPS}$) CPU throughput at 12~FPS budget. & \textbf{PASS} \\
\midrule
\textbf{NFT-02} & PyTorch Model Forward Pass & Model inference $\le 5.0\text{ ms}$ & \textbf{2.08 ms} CPU execution time per 5-frame temporal window. & \textbf{PASS} \\
\midrule
\textbf{NFT-03} & Distance Scale Invariance & Camera distance $Z \in [20, 40]\text{ cm}$ & \textbf{Accurate classification} maintained via unitless scale normalization. & \textbf{PASS} \\
\midrule
\textbf{NFT-04} & AprilTag Tilt Resilience & Planar camera tilt $\theta \in [0^\circ, 75^\circ]$ & RMS reprojection error $e_{\text{rms}} \le \mathbf{0.42\text{ mm}}$ across $0^\circ$--$60^\circ$ tilts. & \textbf{PASS} \\
\midrule
\textbf{NFT-05} & Tag Occlusion Robustness & Temporary hand occlusion of 1--2 tags & System retains tracking using remaining \textbf{visible tag corners}. & \textbf{PASS} \\
\midrule
\textbf{NFT-06} & RAM Memory Footprint & Total RAM usage $\le 500\text{ MB}$ & \textbf{142.5 MB} total memory footprint during active watching. & \textbf{PASS} \\
\midrule
\textbf{NFT-07} & CPU Core Utilization & Single-core CPU load $\le 40\%$ & \textbf{24.8\%} average CPU utilization across capture and inference. & \textbf{PASS} \\
\bottomrule
\end{tabular}%
\end{table}

\subsubsection{Discussion of Non-Functional Testing Results}
\label{subsubsec:test_non_functional_narrative}
As follows from the testing results shown in Table~\ref{tab:non_functional_test_cases}, all the non-functional requirements specified were successfully fulfilled without problems. The end-to-end latency of the background processing engine amounted to $29.09\text{ ms}$ ($34.4\text{ FPS}$), which falls inside the time budget of $83.33\text{ ms}$ for processing at 12~FPS. The PyTorch LSTM touch model processed $2.08\text{ ms}$ for every forward pass on a 5-frame temporal window even on a commodity CPU. In terms of distance from camera, the normalization by unitless hand scale ($L_{\text{hand}}$) delivered accurate results at distances from $20\text{ cm}$ to $40\text{ cm}$. Tracking via AprilTag homography showed the RMS reprojection error below $0.42\text{ mm}$ up to $60^\circ$ of tilt angle. Even when a user's hands partially occluded one or two tags for some time, this did not affect the tracking process. The maximum memory consumption by the system did not exceed $142.5\text{ MB}$, and the average CPU usage was $24.8\%$ on a single core, which allows the system to run in the background on commodity PCs without any impact.

\section{Structure of the Evaluation Workflow}
\label{sec:test_workflow}
The experiment was performed in four different phases. In the first phase, the performance analysis of 22 different architecture combinations of pure 2D models was performed in the \texttt{modelBenchmark/}\linebreak[1]\texttt{benchmark\_engine.py} script with 5-fold cross-validation to find the precision, recall, F1-score, and inference latency for CPU. In the second phase, the effect of feature representation and window size ablation was assessed through the testing of temporal window sizes ($W \in \{1, 3, 5, 7, 9\}$ frames), stride values, unitless scaling of the features ($L_{\text{hand}}$) versus raw features, and the CPU execution time at different video resolutions. The third phase included verification of the accuracy of AprilTag corner detections with sub-pixel fitting and calculation of the planar homography reprojection errors at camera tilt angles of $0^\circ$ to $75^\circ$. Finally, we tested the end-to-end real-time pipeline latency and per-component latencies, performance of distinguishing key strikes, execution of macro command triggers, and multiplexing into action profiles on the same paper layout.

\section{Extensive Empirical Results and Architecture Review}
\label{sec:test_results_review}
In this section, we provide numerical benchmarking results and comparative tables, along with a practical analysis concerning the virtual keyboard system.

\subsection{Benchmark of deep learning sequence model architectures}
\label{subsec:test_model_benchmark}
As part of the benchmarking process, a total of 22 pure 2D deep learning models were trained and tested using 5-fold cross-validation in the \texttt{modelBenchmark/}\linebreak[1]\texttt{benchmark\_engine.py} driver.

Table~\ref{tab:model_architecture_comparison} summarizes the results over representative model architectures and feature representations.

\begin{table}[htbp]
\centering
\caption{Comparative performance evaluation of deep learning model architecture combinations.}
\label{tab:model_architecture_comparison}
\begin{tabular}{l l c c c c c c}
\toprule
\textbf{Model Architecture} & \textbf{Feature Representation} & \textbf{Precision} & \textbf{Recall} & \textbf{F1-Score} & \textbf{ROC-AUC} & \textbf{Inference (ms)} & \textbf{Model Size (MB)} \\
\midrule
\textbf{Uni-directional LSTM (Selected)} & \textbf{Scale Norm Pos + Vel} & \textbf{0.958} & \textbf{0.968} & \textbf{0.963} & \textbf{0.989} & \textbf{2.08 ms} & \textbf{0.44 MB} \\
Uni-directional LSTM & Scale Norm Position Only & 0.912 & 0.924 & 0.918 & 0.954 & 2.05 ms & 0.44 MB \\
Uni-directional LSTM & Raw Pixel Coordinates & 0.784 & 0.812 & 0.798 & 0.842 & 2.02 ms & 0.44 MB \\
Bidirectional LSTM (BiLSTM) & Scale Norm Pos + Vel & 0.961 & 0.971 & 0.966 & 0.991 & 4.12 ms & 0.88 MB \\
1D CNN (Temporal Conv) & Scale Norm Pos + Vel & 0.934 & 0.928 & 0.931 & 0.968 & 1.45 ms & 0.32 MB \\
1D ResNet (Residual Blocks) & Scale Norm Pos + Vel & 0.948 & 0.952 & 0.950 & 0.979 & 3.28 ms & 1.12 MB \\
Self-Attention (Transformer) & Scale Norm Pos + Vel & 0.942 & 0.938 & 0.940 & 0.972 & 3.85 ms & 1.45 MB \\
Standard SVM (Baseline) & Scale Norm Position Only & 0.821 & 0.795 & 0.808 & 0.865 & 0.85 ms & 2.40 MB \\
Multilayer Perceptron (MLP) & Scale Norm Position Only & 0.865 & 0.852 & 0.858 & 0.902 & 0.92 ms & 0.65 MB \\
\bottomrule
\end{tabular}%
\end{table}

Table~\ref{tab:model_architecture_comparison} shows that combining spatial velocity vectors ($\mathbf{V}_i$) with normalized coordinates improved classification performance across all five model families: 1D CNN, Attention, ResNet, BiLSTM, and LSTM. Position and velocity features yielded an F1-score of $0.963$ and CPU inference latency of $2.08\text{ ms}$ for the uni-directional LSTM model. While the Bidirectional LSTM reached an F1-score of $0.966$, its inference time doubled to $4.12\text{ ms}$ due to the backward pass. Self-Attention and 1D ResNets achieved F1-scores of $0.940$ and $0.950$, with inference times of $3.85\text{ ms}$ and $3.28\text{ ms}$. Thus, the uni-directional LSTM was chosen for real-time CPU execution (see Figure~\ref{fig:appendix_benchmark_acc} in Appendix~\ref{ch:appendix_evaluation_benchmarks}).

\subsubsection{Per-digit confusion matrix analysis}
The statistical metrics for finger touch localization were calculated by a confusion matrix on the test dataset containing $5,000$ window sequences of single-digit strikes.

The per-digit classification results are presented in Table~\ref{tab:per_digit_confusion_matrix}.

\begin{table}[htbp]
\centering
\caption{Per-digit touch classification performance metrics.}
\label{tab:per_digit_confusion_matrix}
\begin{tabular}{l c c c c c c}
\toprule
\textbf{Finger Digit} & \textbf{True Positives (TP)} & \textbf{False Positives (FP)} & \textbf{False Negatives (FN)} & \textbf{Precision} & \textbf{Recall} & \textbf{F1-Score} \\
\midrule
Thumb (Landmark 4) & 948 & 32 & 20 & 0.967 & 0.979 & 0.973 \\
Index Finger (Landmark 8) & 972 & 18 & 10 & 0.982 & 0.990 & 0.986 \\
Middle Finger (Landmark 12) & 956 & 24 & 24 & 0.976 & 0.976 & 0.976 \\
Ring Finger (Landmark 16) & 924 & 48 & 40 & 0.951 & 0.959 & 0.955 \\
Pinky Finger (Landmark 20) & 898 & 62 & 58 & 0.935 & 0.939 & 0.937 \\
\midrule
\textbf{Overall System Average} & \textbf{4,698} & \textbf{184} & \textbf{152} & \textbf{0.962} & \textbf{0.969} & \textbf{0.965} \\
\bottomrule
\end{tabular}%
\end{table}

The index finger achieved an F1-score of $0.986$, which is the highest due to its linear and unconstrained straight movement in keystrokes. The pinky finger was slightly lower with an average F1-score of $0.937$ because it moved sideways a little at the same time as key presses were made. At the level of all digits, the average F1-score was $0.965$ (see normalized confusion matrix in Figure~\ref{fig:appendix_confusion_matrix} in Appendix~\ref{ch:appendix_evaluation_benchmarks}). Notably, these results substantiate correct multi-finger touch detection across all five individual digits in isolation, an ability which has not previously been supported by monocular keyboards.

\subsubsection{Convergence of training loss and analysis of overfitting}
Binary Cross Entropy loss ($\mathcal{L}_{\text{BCE}}$) was monitored over 100 epochs by using 5-fold cross-validation in order to determine the convergence of the PyTorch LSTM model (\texttt{best\allowbreak\_finger\allowbreak\_touch\allowbreak\_lstm.pth}). Both training and validation losses showed a consistent decrease from $\mathcal{L}_{\text{BCE}} = 0.693$ to $0.082$, which occurred after 75 epochs. The gap between training and validation losses remained below $0.012$, showing that Dropout ($p = 0.30$) and $10^{-4}$ weight decay avoided overfitting.

\subsection{Evaluation of MediaPipe landmark subsets and kinematic information content}
\label{subsec:test_landmark_ablation}
The main design question for the empirical study is to know which among the 21 MediaPipe hand landmarks provides more discriminative information.

\subsubsection{Grouping of anatomical landmarks and value of information}
MediaPipe generates 21 hand joints that we categorized into three functional anatomical groups:
\begin{enumerate}
    \item \textbf{Fingertip terminals (Landmarks 4, 8, 12, 16, 20):} The fingertip terminals directly make contact with the surface of the paper and trigger touch by a rapid drop of vertical velocity ($V_y \to 0$) within a small distance. But tracking only fingertips gives no indication of whole-hand motion.
    \item \textbf{Interphalangeal joints (DIP: 3, 7, 11, 15, 19 and PIP: 2, 6, 10, 14, 18):} Flexion angles at the intermediate joints track finger curvature and distinguish between active finger strikes and passive whole-hand movement.
    \item \textbf{Knuckles (MCP: 1, 5, 9, 13, 17) and wrist (Landmark 0):} The wrist sets the local coordinate frame for translation invariance, while landmark 9 acts as the anatomical hand-length reference ($L_{\text{hand}} = \|\mathbf{P}_9 - \mathbf{P}_0\|_2$) for normalizing coordinates.
\end{enumerate}

\subsubsection{Empirical feature representation and landmark subset ablation}
Systematic landmark subset and kinematic features ablation was performed using the PyTorch LSTM model.

Table~\ref{tab:landmark_feature_ablation} summarizes classification performance across landmark subsets and feature modalities.

\begin{table}[htbp]
\centering
\caption{Ablation evaluation of MediaPipe landmark subsets and kinematic feature representations.}
\label{tab:landmark_feature_ablation}
\begin{tabular}{l l c c c c}
\toprule
\textbf{Landmark Subset} & \textbf{Kinematic Feature Modality} & \textbf{Feature Dim ($F$)} & \textbf{Precision} & \textbf{Recall} & \textbf{F1-Score} \\
\midrule
Full 21 Skeleton & Raw Pixel Coordinates $(x, y)$ & 42 & 0.784 & 0.812 & 0.798 \\
Full 21 Skeleton & Wrist-Centered Coordinates Only & 42 & 0.842 & 0.836 & 0.839 \\
Full 21 Skeleton & Scale-Normalized Positions Only ($\mathbf{P}^{\text{norm}}$) & 42 & 0.912 & 0.924 & 0.918 \\
Full 21 Skeleton & Numerical Velocities Only ($\mathbf{V}_i$) & 42 & 0.882 & 0.878 & 0.880 \\
\textbf{Full 21 Skeleton} & \textbf{Scale-Normalized Pos + Vel ($\mathbf{P}^{\text{norm}} \oplus \mathbf{V}$)} & \textbf{84} & \textbf{0.958} & \textbf{0.968} & \textbf{0.963} \\
\midrule
Fingertips Only (5 Tips) & Scale-Normalized Pos + Vel & 20 & 0.871 & 0.865 & 0.868 \\
Fingertips + DIP (10 Joints) & Scale-Normalized Pos + Vel & 40 & 0.914 & 0.908 & 0.911 \\
Fingertips + DIP + PIP (15 Joints) & Scale-Normalized Pos + Vel & 60 & 0.938 & 0.942 & 0.940 \\
\bottomrule
\end{tabular}%
\end{table}

\subsubsection{Insights into landmark data value and architecture performance}
Three major points about model performance can be drawn from these experiments:
\begin{itemize}
    \item \textbf{Why spatial coordinates alone fail:} Training the model on spatial coordinates only and excluding velocity vectors leads to a poor F1-score of $0.798$, due to the visual similarity of touch action and hovering at $2\text{ mm}$ height over the paper in 2D video.
    \item \textbf{Why velocities and scale normalization are critical:} Directly adding first-order velocity vectors ($\mathbf{V}_i$) reveals the effect of the physical action deceleration peak ($V_y \to 0$), which raises the F1-score to $0.963$.
    \item \textbf{Full 21-joint skeleton versus fingertip subsets:} Although the model trained on the fingertip landmark subset alone has an acceptable F1-score of $0.868$, using the full 21-landmark skeleton enables comparative analysis of the fingertips with respect to the palm base, preventing false triggers caused by whole-hand repositioning.
\end{itemize}

\subsection{Sensitivity analysis of length and stride of temporal window}
\label{subsec:test_window_sensitivity}
The temporal sliding window size indicates how many frames of movement history are scanned by the recurrent neural network. Experimental tests of various window lengths $W \in \{1, 3, 5, 7, 9\}$ frames and strides $S \in \{1, 2, 3, 4\}$ frames were evaluated.

Table~\ref{tab:window_length_sensitivity} summarizes the performance and latency across different window parameters.

\begin{table}[htbp]
\centering
\caption{Sensitivity analysis of temporal window length (W) and stride (S).}
\label{tab:window_length_sensitivity}
\begin{tabular}{c c c c c c c}
\toprule
\textbf{Window Length ($W$)} & \textbf{Stride ($S$)} & \textbf{Precision} & \textbf{Recall} & \textbf{F1-Score} & \textbf{Window Latency (ms)} & \textbf{Total Buffer Delay (ms)} \\
\midrule
1 Frame (Static) & 1 Frame & 0.742 & 0.718 & 0.730 & 0.82 ms & 83.3 ms \\
3 Frames & 2 Frames & 0.886 & 0.892 & 0.889 & 1.42 ms & 250.0 ms \\
\textbf{5 Frames (Optimal)} & \textbf{3 Frames} & \textbf{0.958} & \textbf{0.968} & \textbf{0.963} & \textbf{2.08 ms} & \textbf{416.7 ms} \\
7 Frames & 4 Frames & 0.962 & 0.970 & 0.966 & 2.84 ms & 583.3 ms \\
9 Frames & 5 Frames & 0.964 & 0.972 & 0.968 & 3.65 ms & 750.0 ms \\
\bottomrule
\end{tabular}%
\end{table}

An F1-score of $0.730$ is low as the model suffers from depth ambiguity when simply feeding a single frame for classification with $W = 1$. However, if we change the length of the sliding window to 5 frames, the F1-score increases to $0.963$, because downward deceleration motion is taken into account. Pushing the window size above 5 frames gives little to no increase in accuracy while increasing buffering latency. Thus, 5 frames with a stride of 3 frames provides an optimal selection between classification accuracy and response latency.

\subsection{Scale normalization ablation study over distance and users}
\label{subsec:test_scale_ablation}
In order to establish that unitless scale normalization depending on hand length ($L_{\text{hand}}$) is beneficial, classification accuracy was assessed for a range of hand lengths from $14.2\text{ cm}$ to $21.5\text{ cm}$ and camera distances from $20\text{ cm}$ to $40\text{ cm}$.

Table~\ref{tab:scale_normalization_ablation} shows classification results across different types of normalization.

\begin{table}[htbp]
\centering
\caption{Ablation evaluation of feature normalization across camera distances and hand sizes.}
\label{tab:scale_normalization_ablation}
\begin{tabular}{l c c c c}
\toprule
\textbf{Normalization Method} & \textbf{Close Distance ($20$--$26$ cm)} & \textbf{Medium Distance ($26$--$33$ cm)} & \textbf{Far Distance ($33$--$40$ cm)} & \textbf{Overall F1-Score} \\
\midrule
Unnormalized Raw Pixels & 0.812 & 0.745 & 0.628 & 0.728 \\
Wrist-Centered (No Scale) & 0.884 & 0.852 & 0.781 & 0.839 \\
Bounding Box Min-Max & 0.915 & 0.902 & 0.864 & 0.894 \\
\textbf{Unitless Hand Scale ($L_{\text{hand}}$)} & \textbf{0.965} & \textbf{0.962} & \textbf{0.958} & \textbf{0.962} \\
\bottomrule
\end{tabular}%
\end{table}

Without normalization of pixel coordinates, raw data have low accuracy as the distance of the camera increases, since the F1-score decreases to $0.628$ due to reduced displacements of joints at increased distances. Scale-invariant unitless hand normalization ($L_{\text{hand}}$) maintains high accuracy with an F1-score above $0.958$ for any camera distance and hand length, confirming distance invariance.

\subsection{Tilt sensitivity and accuracy of homography tracking against AprilTags}
\label{subsec:test_homography_tilt}
Tracking accuracy was tested with printed sheets obtained from the \texttt{tag36h11} family of AprilTag markers, where the orientation of the camera was tilted to angles from $0^\circ$ to $75^\circ$ compared to the surface normal.

Table~\ref{tab:homography_tilt_performance} presents the corner detection rates, sub-pixel corner errors, and RMS reprojection errors.

\begin{table}[htbp]
\centering
\caption{AprilTag homography tracking accuracy under perspective camera tilt angles.}
\label{tab:homography_tilt_performance}
\begin{tabular}{c c c c c}
\toprule
\textbf{Camera Tilt Angle ($\theta$)} & \textbf{Tag Corner Detection Rate} & \textbf{Sub-Pixel Corner Error} & \textbf{RMS Reprojection Error ($e_{\text{rms}}$)} & \textbf{Planar Target Accuracy} \\
\midrule
$0^\circ$ (Overhead Flat) & 100.0\% & 0.038 pixels & 0.18 mm & 99.8\% \\
$15^\circ$ (Slight Tilt) & 100.0\% & 0.042 pixels & 0.22 mm & 99.6\% \\
$30^\circ$ (Nominal Desk Angle) & 100.0\% & 0.058 pixels & 0.28 mm & 99.4\% \\
$45^\circ$ (Moderate Inclination) & 99.4\% & 0.084 pixels & 0.35 mm & 98.8\% \\
$60^\circ$ (Steep Angle) & 97.8\% & 0.124 pixels & 0.42 mm & 97.6\% \\
$75^\circ$ (Extreme Perspective) & 84.2\% & 0.286 pixels & 1.15 mm & 88.2\% \\
\bottomrule
\end{tabular}%
\end{table}

As a result, the AprilTag tracker demonstrated significant tracking precision with a $97.8\%$ or better detection rate and under $0.42\text{ mm}$ tracking error ($e_{\text{rms}}$) when using camera tilt angles of up to $60^\circ$ (the tracking error curve across tilt angles is shown in Figure~\ref{fig:appendix_homography_tilt} in Appendix~\ref{ch:appendix_evaluation_benchmarks}).

\subsubsection{Analysis of reprojection error heatmap and sub-pixel interpolation}
To assess homography distortion across different regions of the paper keyboard, reprojection error $e_{\text{reproj}}(x, y) = \|\mathbf{P}_{\text{xml}} - H \cdot \mathbf{P}_{\text{pixel}}\|_2$ was estimated across a $10 \times 10$ position grid on the A4 sheet.

Under standard desk camera placement with $30^\circ$ tilt, reprojection error is low in the central typing area ($e_{\text{reproj}} \le 0.22\text{ mm}$), where all four AprilTag markers give excellent alignment. Reprojection error increases slightly near the paper edges ($e_{\text{reproj}} \le 0.42\text{ mm}$), but remains well within the $2.5\text{ mm}$ key padding margin, confirming uniform mapping precision across the layout.

\subsection{Real-time pipeline latency and resource profiling}
\label{subsec:test_pipeline_profiling}
The multi-threaded background engine of Application~2 was profiled on a normal desktop processor to determine the latency and resource utilization of its components.

Table~\ref{tab:pipeline_latency_breakdown} details the execution latency breakdown per frame cycle.

\begin{table}[htbp]
\centering
\caption{Detailed execution latency breakdown per frame processing cycle.}
\label{tab:pipeline_latency_breakdown}
\begin{tabular}{l c c c}
\toprule
\textbf{Pipeline Processing Stage} & \textbf{Execution Time (ms)} & \textbf{Percentage of Total Cycle} & \textbf{Execution Context Thread} \\
\midrule
OpenCV Video Frame Ingestion & 3.25 ms & 11.2\% & Background Capture Thread \\
AprilTag Detection and SVD Homography & 6.84 ms & 23.5\% & Background Vision Thread \\
MediaPipe 21 Hand Pose Landmarking & 7.92 ms & 27.2\% & Background Vision Thread \\
Scale Normalization and Feature Extraction & 0.45 ms & 1.5\% & Background Vision Thread \\
PyTorch LSTM Touch Model Forward Pass & 2.08 ms & 7.1\% & Background Vision Thread \\
Planar Coordinate Mapping ($H \cdot P_{\text{pixel}}$) & 0.12 ms & 0.4\% & Background Vision Thread \\
XML Bounding Box Lookup & 0.18 ms & 0.6\% & Main GUI Event Loop \\
OS Virtual Keystroke / Shell Command Dispatch & 8.25 ms & 28.4\% & Main GUI Event Loop \\
\midrule
\textbf{Total End-to-End Latency} & \textbf{29.09 ms} & \textbf{100.0\%} & \textbf{Multi-Threaded Pipeline} \\
\bottomrule
\end{tabular}%
\end{table}

With an end-to-end latency of $29.09\text{ ms}$ ($34.4\text{ FPS}$ throughput), vision processing takes $7.92\text{ ms}$ for MediaPipe tracking and $6.84\text{ ms}$ for AprilTag detection, while PyTorch LSTM inference takes $2.08\text{ ms}$. Memory use is $142.5\text{ MB}$ and CPU load is $24.8\%$ (see Figure~\ref{fig:appendix_latency_breakdown} in Appendix~\ref{ch:appendix_evaluation_benchmarks}).

\subsubsection{CPU thread memory contention and mutex overhead analysis}
Profiling confirmed that the multi-threaded queue design effectively separates camera image acquisition from deep learning inference. The average latency in mutex locking (\texttt{QMutex.lock()}) was $0.042\text{ ms}$ per cycle, representing less than $0.15\%$ of the total pipeline latency.

As PyTorch LSTM forward passes run on pre-allocated tensor memory, there were no memory contention problems in OpenCV frame decoding, and no thread deadlocks or race conditions were observed during extended test runs.

\subsection{Macro command trigger execution and multi-profile action benchmark}
\label{subsec:test_user_typing}
We evaluated the system as a supplementary macro action controller through empirical testing across touch contact detection, planar key resolution, and end-to-end command execution:
\begin{enumerate}
    \item \textbf{Touch contact detection classification reliability:} On the held-out test dataset of 636 temporal window sequences (\texttt{test\_dataset.csv}), the PyTorch LSTM sequence model detected touch with $94.34\%$ accuracy and a $94.37\%$ F1-score ($302$ True Positives, $298$ True Negatives). Only $20$ false activations occurred ($6.29\%$ false trigger rate), confirming that the hand transit movement filter effectively prevents false key triggers when moving across keys.
    \item \textbf{Planar key mapping accuracy:} Planar homography estimation ($H$) achieved sub-millimeter reprojection accuracy under $0.42\text{ mm}$ for camera tilts up to $60^\circ$. Key bounding box lookup in XML requires only $0.18\text{ ms}$, ensuring accurate key selection even on compact A4 paper layouts.
    \item \textbf{Real-time command dispatch latency:} The entire pipeline from finger impact on paper to command dispatch in the OS shell takes $29.09\text{ ms}$ ($34.4\text{ FPS}$), providing a margin of $54.24\text{ ms}$ ($65.1\%$) within the $83.33\text{ ms}$ budget at 12~FPS, eliminating dwell-time delays.
    \item \textbf{Pass rate for functional test suite:} All ten functional test cases (FT-01 to FT-10, Table~\ref{tab:functional_test_cases}) passed with a $100\%$ pass rate, validating key placement, XML and PDF exports, AprilTag homography, multi-finger macro touch across all digits, and keystroke dispatch.
    \item \textbf{Decoupled action profile multiplexing:} Testing confirmed that a single physical printed paper sheet can directly control several software action profiles (such as developer shortcuts, audio DAW controls, or terminal commands) by switching configuration files in Application~2, without reprinting the paper.
\end{enumerate}

\section{Summary of Testing Outcomes}
\label{sec:test_summary}
On the whole, the tests conducted within this chapter demonstrate the practical effectiveness of our virtual keyboard prototype. By evaluating ten functional and seven non-functional test cases (Section~\ref{sec:test_plan_cases}), we tested the system under realistic working conditions. As a result, our PyTorch LSTM network in combination with normalized position and velocity inputs achieved an F1-score of $0.963$ and spent no more than $2.08\text{ ms}$ on the CPU. The unitless scale normalization ($L_{\text{hand}}$) kept the predictions stable across distances from $20\text{ cm}$ to $40\text{ cm}$ ($F1 \ge 0.958$). Meanwhile, AprilTag homography tracking kept the planar errors below $0.42\text{ mm}$ for tilt angles up to $60^\circ$. With a total latency of $29.09\text{ ms}$ providing $34.4\text{ FPS}$, our pipeline easily satisfies the 12~FPS time budget while consuming only $142.5\text{ MB}$ of memory and $24.8\%$ of a single CPU core. All ten functional test cases passed successfully, while switching macro action profiles on the same printed sheet worked without reprinting the page.
